% !TeX root = ../papers/nc-p3.tex
\section{量子境界の五中心局所方程式と対角Ore枝との全交叉}
\label{sec:quantum-ore-incidence}
\begin{theorem}\label{thm:quantum-ore-incidence}
定理\ref{thm:quantum-component-intersection}の量子点の近くで、二隣接階数6のテンソルのうち
四次中心が五次元以上となる局所スキームは滑らかで、ちょうど十五次元成分に一致する。
さらに、基底の一径数部分群による射影極限がこの量子点である任意の隣接階数6テンソルは、
この成分に属し、それ自身が全次数で標準Koszul左右AS正則である。
したがって量子テンソルへのこの形の退化を持つ全ての窓で、成分所属は極限の量子同型類だけで決まる。
対角量子Ore表示 $A(q,f)$ が標準AS正則であるとする。
そのindexed類がこの成分の閉包にある必要十分条件は、背景 $A(q,0)$ が同定理の唯一の量子類であることである。
また定理\ref{thm:quantum-tangent-finite}の二つの九係数表示では $\zeta^2=-1$ の全係数がこの閉包に属し、
$\zeta=-1$ の表示は係数の特殊値を含めて全て外にある。有限四辺はここでは仮定しない。
\end{theorem}
\begin{proof}[階数条件の局所成分]
量子点で経路長4,5の経路加群は隣接階数スキームの近傍で局所自由である。
実際節\ref{sec:unrestricted-doublet}の局所Artin環上の平坦性を完備化の各商に適用し、
有限表示加群の完備化と忠実平坦性を用いればよい。
中心の交換行列の核が五次元以上となる小行列式スキームを $Z_5$ とする。
量子点で核はちょうど五次元であり、$Z_5$ 上の一次変形ではその基底が持ち上がる。

二隣接切断のSchur微分の全階数はこの点で少なくとも177である。
\path{computations/verify_quantum_ore_incidence.py} は四乗根を法97で22として、
量子交代テンソルの両切断の左右の十次元核で摂動を挟み、二百行の階数177を得る。
計算では三辺を1にしたゲージを使う。巡回標準化のための四基底変更はこの階数を変えない。
よってアフィン接空間次元は高々79である。三つの三角積の量子変形に制限すると、
$Z_5$ の一次条件の核は零である。実際体積単項式を中心にする前節のゲージは無限小でも選べる。
四つの純四乗の持上げを多重次数ごとに射影すると同じ純単項式の非零倍に限られる。
従って一階の交代指標の変化は全 $4e_i-\mathbf1$ と $L$ の各元との組で零になる。
四ベクトルは $L\otimes\mathbb Q$ を張るので三方向は全て消える。
従って $\dim T Z_5\le76$ である。他方既約な十五次元成分のアフィン閉包は次元76で、
一般点の五中心から $Z_5$ に含まれ量子点を通る。よって局所次元も接次元も76で、
局所環は正則、成分は一意である。
\end{proof}

\begin{proof}[量子点へ流入する全ての窓]
量子点を固定する連結射影安定化群は四頂点に共通な対角トーラスである。
実際、四三角積は全て1と異なるので四頂点の点配置は六本の座標直線で、その四三重点は内在的である。
連結群は各頂点で対角に作用する。六つの量子関係の両項の比を保つには、
位置の異なる対角行列は射影的に同じでなければならない。
このトーラスは三つの量子パラメータを変えたテンソルの直線も全て固定する。
アフィン空間ではその直線上の同じ倍率を除いて作用させる。従って
$T Z/T Z_5$ は自明指標だけからなる。実際、上の三つの量子接方向が商で独立で、
$\dim T Z\le79$ と $\dim T Z_5=76$ から商の全体を張る。

任意の一径数作用で量子点 $w$ を固定するものを取る。
$Z_5$ の完備局所イデアルはこの作用の重み零の三元で生成されることを確認する。
$Z$ の最大イデアルの余接空間の固有基底を選ぶ。$Z_5$ への制限の核の三ベクトルは重み零なので、
それらを $Z_5$ のイデアルから重み零で持ち上げ、残りを固有元で持ち上げる。
各有限jetでトーラス表現の重み分解が完全であるから、これらの選択は完備化で両立してできる。
全持上げは完備局所環を位相的に生成する。最後の76元は滑らかな $Z_5$ の正則形式座標なので、
最初の三元で割った後、76変数形式冪級数環からの全射とその $Z_5$ への写像の合成は同型である。
従って三元だけでイデアルを生成する。これは $Z$ の他の局所成分を排除する主張ではない。

いま $W$ の軌道弧 $t\mapsto\lambda(t)W$ の極限を $w$ とする。射影極限では
最小重みの全体倍率も除いて、非零なアフィン極限を取る。作用はこの極限を固定し、
余接商への作用は上の連結安定化群を通じて自明である。
三つの重み零生成元の $k[[t]]$ への引戻しは $t\mapsto st$ に不変だから定数で、定数項は零である。
よってこの弧は $Z_5$ の一意な局所成分に入り、四基底変更で不変な閉包には $W$ 自身が入る。
さらに特殊点が標準AS正則なので、付値環上の命題\ref{prop:unrestricted-good-reduction}で
分数体側も全次数で正則であり、軌道上の $W$ も正則となる。
逆に他の量子点へこの形で退化する窓が十五次元成分にあれば、基底変更不変な閉包にその量子点も属し、前節に反する。
\end{proof}

\begin{proof}[九係数表示の中心]
純三文字の巡回量子係数を $\zeta$、$\zeta^2=-1$ とする。混合式を
$X_n x_a-x_a X_{n+1}=f_{n,a}$ と書く。各項は $a$ を含まない $x_b^s x_c^{2-s}$ であり、
係数の位置倍率は $\zeta^{n(s+1)}$ である。この表示の厳密四移動はFrobenius標準である。
これを確認するため、双対の量子外積への重み付き随伴次数とその最高次基底でtraceを規格化する。
PBW基底で各次数の双対と対合は全係数空間上で自由、完全であるから、
局所Nakayama行列は係数に正則であり、対角基底倍率に共変である。
$X$ だけの倍率により各 $C_{a,s}$ は重み1にできる。従って基点と異なり得るのは
$X$ と純文字の間の非対角成分だけで、そこでも係数は一次でなければならない。
純文字の三つの独立な倍率で比べると、$C_{a,s}$ の純部分の指標は
$e_a-se_b-(2-s)e_c$（全指標の符号は作用の向きによる）であり、
一文字と $X$ の間に必要な $\pm e_j$ に等しくない。従って補正はない。
背景では全行積が1なので標準四移動は通常の反復である。

純三文字の量子Laurent環に一つの可換文字 $Y$ を加えた窓に、
$X_n=Y+\sum_a g_{n,a}$ と置いて埋め込む。$g_{n,a}$ は三つのLaurent単項式
$x^{\alpha_s}$、$\alpha_s=-e_a+s e_b+(2-s)e_c$ の線型結合である。
$x_j x^{\alpha_s}=\chi_j x^{\alpha_s}x_j$ とすれば
$\chi_b=\chi_c=\zeta^{-s-1}$、$\chi_a=\zeta^{2s-2}$ である。
従って位置倍率を $\chi_b^{-n}$ とすると $b,c$ で混合差は零となり、
$a$ では $1-\chi_a/\chi_b\ne0$ であるから所定の係数を作れる。
PBW順序で $X$ についてモニックな置換なのでこの表現は忠実である。

一つの極だけを残した $U_{n,a}=Y+g_{n,a}$ と
$P_{n,a}=U_{n,a}U_{n+1,a}U_{n+2,a}U_{n+3,a}$ を作る。$P_{n,a}$ はその一極窓の中心切断である。
$b,c$ とは $U$ 自身が交換し、$a$ を四回移した誤差は各 $s$ で
$\sum_{r=0}^3\zeta^{r(s+1)}=0$ を因子に持つ。実際、誤差の純二文字は $b,c$ だけなので
残る三つの $U$ を全て同じ側へ移せる。また $U$ とは四周期で交換する。
Laurent化すると $P_{n,a}$ は $Y$ とも交換するので $P_{n,a}=P_{n+1,a}$、
各純文字と通常の意味で可換である。以後 $P_a$ と書く。
$Y^{4-r}$ の係数に現れる純指数は $a$ 方向で $-r$、他の二方向で非負である。
三文字のLaurent中心の指数は三成分が法4で等しいので、$1\le r\le4$ では
他の二指数がともに $4-r$ 以上となる。

全三極に対し
\[
             z=Y^4+\sum_{a=0}^2(P_a-Y^4)
\]
はLaurent窓の中心切断であり、実際に元の窓に属する。
固定した $a$ 以外の純文字を可逆にすると、$U_{n,a}=X_n-\sum_{b\ne a}g_{n,b}$ はこの局所表示に入る。
$P_a$ も入る。$b\ne a$ の $P_b-Y^4$ の $Y^j$ の係数は $x_a^j$ で割れる。
$Y=U_{n,a}-g_{n,a}$ を代入すると $g_{n,a}$ の $x_a$ の極は高々一位である。
厳密には $Y^j$ の代わりに、順序付き $U$ 積の最高 $Y$ 項を順に消す。
長さ $d$ の $U$ 積の $Y^h$ 係数の付値は $h-d$ 以上なので、残る $Y^h$ の係数は
常に $x_a$ で少なくとも $h$ 回割れ、全ての $U$ 係数が整である。
これは各 $a$ について成立する。各局所表示の $X$ 順序基底は同じなので、
その係数の交叉は元の量子多項式環であり $z$ が整となる。
同様に $x_0x_1x_2Y$ では三極が一度ずつ消える。
これら二式と純四乗三式は $Y$ 次数と純多重次数により独立な五つの中心四次切断である。
前掲コードは三極を同時に持つ例で、四位置の $X$ 順序基底への書戻しも検算する。
\end{proof}

\begin{proof}[枝全体への適用]
九係数を同時に $t$ 倍する族は $t\ne0$ で $X$ の倍率変更により一定同型類で、
$t=0$ で量子背景になる。原始四乗根の背景は前節の唯一の類である。
全族に五中心があるため、量子点の $Z_5$ の一意な局所成分の結果が族全体に適用される。
一方 $\zeta=-1$ で閉包に入る点があれば同じ基底変更による特殊化でその量子背景も入るが、
その三角積の群は位数2であり前節に反する。

任意の正則な対角量子Ore表示でも、補題\ref{lem:quantum-ore-background-rank}の縮尺族により
背景の所属が必要である。逆に背景が指定量子類なら、第四文字との三つの交換係数を対角twistで1にする。
残る三辺は三角積であり、純三文字を付け替えて巡回順に同じ原始四乗根となる。
混合補正は位置依存の純二次式となり、式\eqref{eq:ternary-shift-derivation}を満たす。
定理\ref{thm:quantum-tangent-finite}の多重次数分解では、非負の $e_a$ 次数の補正は
$\mathbf1,\chi$ が独立なので位置別の $X$ への一次式の加算で消える。
残る次数は上記九係数だけで、その位相も同じ式で定まる。
この除去計算は辺の有限性を使わない。従って標準化したindexed窓は原始四乗根の九係数表示に入り、
十分性も従う。
\end{proof}
